\section{Experimental Setup}
\label{sec:experimental}

The experimental protocol uses fixed chronological splits and a consistent preprocessing and model-evaluation configuration. Training uses the earliest chronological block, validation begins on 2024-01-01, and the untouched test interval begins on 2024-12-28. No random row shuffle is used. Metrics and figures are computed from the cleaned dataset using the same feature-generation, model-inference, Stage-1 clipping, and post-processing definitions described in this paper. The test should be read as an offline target-weather-input evaluation: it does not include archived issue-time forecasts, forecast lead times, or weather-source forecast error.

\begin{table}[!t]
\caption{Row-count reconciliation for the final test evaluation.}
\label{tab:evaluation_filter_summary}
\centering
\small
\begin{tabular}{l r}
\toprule
Evaluation step & Rows \\
\midrule
Raw test interval starting 2024-12-28 & 106,272 \\
Rows with complete proposed-model features & 106,272 \\
Rows with complete sensor-model features & 106,272 \\
Daylight mask, proposed path & 52,731 \\
Common mask with 5-min persistence & 52,731 \\
Common mask with previous-day persistence & 52,598 \\
\bottomrule
\end{tabular}
\end{table}


Table~\ref{tab:evaluation_filter_summary} reconciles the row counts that appear in the paper. The test interval contains 106,272 raw 5-min rows. The headline daylight evaluation uses rows with $\Gcs>10$ W/m$^2$, giving 52,731 proposed-path samples. The Stage-1 $\kt$ training filter also records a separate solar-elevation threshold in the metadata; this threshold is expressed in degrees and is not the same thing as the final daylight mask used for the headline power metrics.

\begin{table*}[!t]
\caption{Evaluation masks and row-count predicates used for the reported experiment.}
\label{tab:evaluation_masks}
\centering
\scriptsize
\setlength{\tabcolsep}{4pt}
\begin{tabularx}{\textwidth}{@{}>{\raggedright\arraybackslash}p{0.20\textwidth}>{\raggedright\arraybackslash}X>{\raggedleft\arraybackslash}p{0.12\textwidth}@{}}
\toprule
Use & Predicate & Rows \\
\midrule
Retained 5-min grid & Sorted unique IST timestamps after data cleaning. Machine-fault, hardware-error, shutdown, and invalid-sensor periods are omitted before this retained grid is modeled. & 525,888 \\
Stage-2 training block & $t<\mathrm{2024\mbox{-}01\mbox{-}01}$ and all active Stage-2 features available. The first 12 raw training rows are unavailable because 1-h lag features require previous samples. & 315,348 \\
Validation block & $\mathrm{2024\mbox{-}01\mbox{-}01}\leq t<\mathrm{2024\mbox{-}12\mbox{-}28}$ with the active feature list complete. & 104,256 \\
All test block & $t\geq\mathrm{2024\mbox{-}12\mbox{-}28}$ with the active feature list complete. & 106,272 \\
Stage-1 training mask & Stage-2 training block plus $\alpha(t)>5^\circ$; the target is the clipped $\kt$ definition in Eq.~\eqref{eq:kt}. & 148,269 \\
Stage-1 test mask & All test block plus $\alpha(t)>5^\circ$ for Stage-1 transmissivity diagnostics. & 49,923 \\
Headline daylight test & All test block plus $\Gcs(t)>10$ W/m$^2$ for power-forecast reporting. & 52,731 \\
Previous-day common mask & Headline daylight test plus available $P(t-24~\mathrm{h})$. & 52,598 \\
\bottomrule
\end{tabularx}
\end{table*}


\begin{figure}[!t]
\centering
\includegraphics[width=\columnwidth]{figures/diagrams/evaluation_filter_flow.png}
\caption{Evaluation filter flow and row-count reconciliation for the final test evaluation.}
\label{fig:evaluation_filter_flow}
\end{figure}

\subsection{Metrics}

For observations $y_i$ and predictions $\hat{y}_i$, the mean absolute error is
\begin{equation}
\mathrm{MAE}=\frac{1}{N}\sum_{i=1}^{N}|y_i-\hat{y}_i|.
\label{eq:mae}
\end{equation}
The root mean square error is
\begin{equation}
\mathrm{RMSE}=\sqrt{\frac{1}{N}\sum_{i=1}^{N}(y_i-\hat{y}_i)^2}.
\label{eq:rmse}
\end{equation}
The normalized MAE and RMSE are reported relative to plant rating:
\begin{equation}
\mathrm{nMAE}=\frac{\mathrm{MAE}}{\Prated}\times100\%, \qquad
\mathrm{nRMSE}=\frac{\mathrm{RMSE}}{\Prated}\times100\%.
\label{eq:normalised_errors}
\end{equation}
The coefficient of determination is
\begin{equation}
R^2=1-\frac{\sum_{i=1}^{N}(y_i-\hat{y}_i)^2}
{\sum_{i=1}^{N}(y_i-\bar{y})^2}.
\label{eq:r2}
\end{equation}
The Nash-Sutcliffe efficiency (NSE) is computed with the same residual-to-variance ratio:
\begin{equation}
\mathrm{NSE}=1-\frac{\sum_{i=1}^{N}(y_i-\hat{y}_i)^2}
{\sum_{i=1}^{N}(y_i-\bar{y})^2}.
\label{eq:nse}
\end{equation}
Forecast skill against a baseline $b$ is
\begin{equation}
\mathrm{Skill}_b=1-\frac{\mathrm{RMSE}_{\mathrm{model}}}
{\mathrm{RMSE}_{b}}.
\label{eq:skill}
\end{equation}
Skill is reported only when the model and baseline are evaluated on the same rows.

\subsection{Persistence Baselines}

Two persistence baselines are reported because they answer different operational questions. The 5-min last-value baseline is
\begin{equation}
\hat{y}^{\mathrm{LV}}(t)=y(t-5~\mathrm{min}).
\label{eq:last_value_persistence}
\end{equation}
This is a strong nowcast because it uses live plant output immediately before the target timestamp. The previous-day same-time baseline is
\begin{equation}
\hat{y}^{\mathrm{PD}}(t)=y(t-24~\mathrm{h}).
\label{eq:previous_day_persistence}
\end{equation}
This is a simpler schedule-oriented forecast baseline. The proposed model is compared against both on common daylight masks. Reporting both prevents an inflated skill claim: the proposed weather-input path improves over previous-day persistence, but it does not beat the 5-min last-value nowcast under the current test protocol.

\subsection{Ablations}

Two ablations are evaluated with the same feature pipeline. The first is a single-stage XGBoost model that predicts power directly from weather, time, solar, and cloud features without the learned $\kt$ intermediate. The second preserves the two-stage cascade but removes low-, mid-, and high-cloud channels, leaving only total cloud cover and its derived lag/rolling features. Both ablations reuse the chronological split and tuned parameter family so the comparison isolates model structure rather than data leakage or a different test interval.

\subsection{Implementation Notes}

The implementation supports GPU-accelerated XGBoost for faster experimentation, with CPU execution as an equivalent fallback. Hardware acceleration is not a scientific contribution of the paper; it affects only training convenience. The deployment inference path accepts a weather forecast CSV, interpolates and engineers features, loads the saved models, clips the prediction envelope, and writes prediction outputs. Fig.~\ref{fig:inference} summarizes this path; the reported test metrics still come from the offline held-out target-weather-input protocol.

\begin{table*}[!t]
\caption{Reported model configuration used for the experimental evaluation.}
\label{tab:model_configuration}
\centering
\scriptsize
\setlength{\tabcolsep}{3pt}
\begin{tabularx}{\textwidth}{@{}l>{\raggedright\arraybackslash}X>{\raggedright\arraybackslash}Xcccc@{}}
\toprule
Model & Role & Rows used for fit & Features & Estimators & Depth & Learning rate \\
\midrule
Stage 1 XGBoost & $\kt$ estimation & 148,269 daylight train rows & 36 & 732 & 6 & 0.0218 \\
Stage 2 XGBoost & Proposed power forecast & 315,348 train rows & 41 & 491 & 3 & 0.0459 \\
Sensor XGBoost & Diagnostic upper bound & 315,348 train rows & 41 & 491 & 3 & 0.0651 \\
\bottomrule
\end{tabularx}
\end{table*}


Table~\ref{tab:model_configuration} records the model family used in the reported experiments. Stage 1 is trained only on normal daylight rows because $\kt$ is unstable near night and very-low-sun periods. Stage 2 uses the full training block after feature cleaning so the final power model sees night, transition, and daylight behavior. The sensor model has the same broad feature count as Stage 2, but replaces the estimated irradiance family with measured irradiance-derived features. The ablation models in Table~\ref{tab:ablation_summary} use the same chronological split as the proposed model, rather than a separate or stale evaluation interval.

\begin{figure}[!t]
\centering
\includegraphics[width=\columnwidth]{figures/diagrams/diagram_inference.png}
\caption{Deployment inference path. The same feature-generation logic is reused when weather forecast data are supplied through a CSV file before prediction issue time.}
\label{fig:inference}
\end{figure}
